\paragraph*{Experimental scales.}
The two control parameters are
$\beta=2\pi L I_c/\Phi_0$ and $E_J/\hbar\omega$, related by
$\beta=2\pi(Z/R_Q)(E_J/\hbar\omega)$, with $E_J/h\simeq0.5\,{\rm GHz}\times I_c[{\rm nA}]$.
Three classes of $\pi$ elements span the phase diagram
[Fig.~\ref{fig:overview}(b)]:
(i) SFS junctions ($I_c\sim1$--$100\,\mu$A) reach $\beta>1$ with geometric
inductances of a few to a few hundred pH and $E_J/\hbar\omega\sim10^2$--$10^4$.
This is the classical limit, where the half-flux state is fully polarized,
consistent with the spontaneous currents observed in SFS
and $d$-wave $\pi$ loops~\cite{Tsuei2000,Ryazanov2001,Bauer2004}.
(ii) Quantum-dot $\pi$ junctions in the doublet regime
($I_c^\pi\sim0.1$--$10$\,nA) require $L\sim0.03$--$3\,\mu$H, i.e.\ superinductances
(Josephson-junction arrays, granular Al) as in fluxonium circuits~\cite{Manucharyan2009,Grunhaupt2019}; this is the $\pi$-qubit regime of Refs.~\cite{Ioffe1999,Yamashita2005,Feofanov2010}. These give
$E_J/\hbar\omega\sim1$: the ground state is the
cat state $(|{\circlearrowright}\rangle+|{\circlearrowleft}\rangle)/\sqrt2$ with
$\delta\sim0.1\hbar\omega$, and the transition appears as a dip in the
$|0\rangle\to|1\rangle$ frequency and a jump in the transition matrix element
$\langle0|\Phi|1\rangle$, accessible in two-tone spectroscopy.
(iii) Multichannel or gate-tunable $\pi$ junctions with
$I_c\sim50$--$200$\,nA in a high-kinetic-inductance loop ($L\sim2$--$10$\,nH,
$\omega/2\pi\sim2$--$5$\,GHz) give $\beta\sim1$--$3$ and $E_J/\hbar\omega\sim10$--$30$.
This is the regime of Figs.~\ref{fig:halfflux} and~\ref{fig:udot}: a well-defined
current-carrying state whose residual tunnelling $\delta/\hbar\omega\sim10^{-4}$--$10^{-8}$
remains a measurable macroscopic-quantum-tunnelling signature.
In the externally biased realization, the half-flux bias must be stable to
$|\delta\Phi|\ll\delta/2I^\star$, i.e.\ $\lesssim10^{-6}\Phi_0$ for the parameters above.
The interaction-driven realization has no such requirement.
